%% ============================================================
%% cap06.tex — Capítulo 6
%% Pruebas, calidad y refactorización
%% Sprints 10 y 11
%%
%% Autor: Roberto Albeiro Pava-Díaz, Ph.D.
%% ORCID: 0000-0003-0440-892X
%% Universidad Distrital Francisco José de Caldas
%% Versión 2026 — Fase 2 completa
%% ============================================================

\chapter{Pruebas, calidad y refactorización}
\label{cap:pruebas}

\noindent
Un sistema de información que implementa estatutos universitarios debe ser no solo funcional,
sino correcto: las invariantes del dominio ---que una Escuela no puede existir sin Facultad,
que un Estudiante matriculado en un programa tiene exactamente un plan de estudios activo,
que los tipos de docente son los cinco definidos en el artículo~8 del \acuerdo--- no son
sugerencias de diseño; son restricciones con consecuencias normativas. La diferencia entre
un sistema que «parece funcionar» y uno que \emph{puede demostrarse} que funciona reside,
casi por completo, en la existencia de una batería de pruebas automatizadas que se ejecuta
en cada integración del código.

Queda abierta, sin embargo, la pregunta de cuántas pruebas son suficientes. La cobertura del
100~\% de líneas no garantiza la ausencia de errores lógicos; tampoco el cumplimiento de todas
las aserciones de dominio garantiza que el sistema responda correctamente ante interacciones
inesperadas del usuario. Esa tensión irreducible es una de las razones por las cuales la
prueba automatizada se trata, en este capítulo, no como un fin en sí mismo, sino como parte
de un ciclo más amplio que incluye métricas de calidad, refactorización continua y
documentación acuciosa.\footnote{En la Facultad de Ingeniería de la \ud, el autor ha
observado ---entre 2019 y 2025, en las asignaturas Programación Orientada a Objetos y
Patrones de Diseño de Software--- que los estudiantes que escriben pruebas \emph{antes}
de implementar la funcionalidad (TDD) producen código con menor acoplamiento y mayor
cohesión, aunque la diferencia se percibe con claridad solo a partir del tercer sprint
del proyecto.}

Este capítulo cubre los Sprints~10 y~11. El primero se orienta a la construcción de la
batería de pruebas unitarias para el dominio; el segundo, a la refactorización fundamentada
en los principios SOLID \parencite{Martin2017} y en el catálogo de técnicas de Fowler
\parencite{Fowler2018}. Ambos sprints se apoyan en las clases del paquete
\texttt{co.edu.udistrital.estatuto.modelo} implementadas en el Capítulo~\ref{cap:java}
y en los \textit{structs} del módulo \texttt{modelo} de la implementación en \rust\
del Capítulo~\ref{cap:rust}.


%% ─────────────────────────────────────────────
\section{Pruebas unitarias en \java\ (JUnit~5)}
\label{sec:pruebas-java}
%% ─────────────────────────────────────────────

\subsection{Arquitectura de JUnit~5}
\label{subsec:junit5-arquitectura}

JUnit~5 no es una biblioteca monolítica, sino una plataforma compuesta por tres módulos
con responsabilidades bien delimitadas, lo que permite su extensión sin modificar el núcleo.
La tabla~\ref{tab:junit5-arquitectura} resume los componentes y sus funciones.

\begingroup
\small
\begin{longtable}{@{} L{3.5cm} L{4cm} L{6.5cm} @{} }
	% --- Título y encabezado de la primera página ---
	\caption{Componentes de la plataforma JUnit~5.} \label{tab:junit5-arquitectura} \\
	\toprule
	\textbf{Módulo} & \textbf{Función principal} & \textbf{Notas} \\
	\midrule
	\endfirsthead
	
	% --- Encabezado para las páginas siguientes ---
	\multicolumn{3}{c}{{\bfseries \tablename\ \thetable{} -- Continuación}} \\
	\toprule
	\textbf{Módulo} & \textbf{Función principal} & \textbf{Notas} \\
	\midrule
	\endhead
	
	% --- Pie de tabla para páginas intermedias ---
	\midrule
	\multicolumn{3}{r}{{Continúa en la siguiente página...}} \\
	\bottomrule
	\endfoot
	
	% --- Pie de tabla final ---
	\bottomrule
	\endlastfoot
	
	% --- Datos de la tabla ---
	\texttt{JUnit Platform} &
	Motor de descubrimiento y ejecución &
	Define la API \texttt{TestEngine}; soporta múltiples \textit{frameworks} sobre una sola plataforma \\
	\midrule
	
	\texttt{JUnit Jupiter} &
	API de escritura de pruebas y extensiones &
	Introduce anotaciones modernas como \texttt{@Test}, \texttt{@BeforeEach}, \texttt{@Nested} y \texttt{@ParameterizedTest} \\
	\midrule
	
	\texttt{JUnit Vintage} &
	Compatibilidad con versiones anteriores &
	Permite ejecutar suites heredadas (JUnit~3 y 4) junto con pruebas Jupiter sin migración forzosa \\
\end{longtable}
\endgroup

Las anotaciones más relevantes para el proyecto del estatuto se listan a continuación.
No es descartable que en versiones subsecuentes de JUnit~5 se incorporen nuevas anotaciones
---al momento de redactar este capítulo (marzo de 2026), la versión estable es JUnit~5.10---,
pero las aquí descritas conforman el repertorio estable desde la versión~5.5:

\begin{itemize}
  \item \texttt{@Test}: marca un método como caso de prueba individual.
  \item \texttt{@BeforeEach} / \texttt{@AfterEach}: ejecutan código de preparación y limpieza
        antes y después de \emph{cada} prueba; análogos a \texttt{setUp()} y \texttt{tearDown()}
        de JUnit~4, pero sin herencia obligatoria.
  \item \texttt{@BeforeAll} / \texttt{@AfterAll}: ejecutados una sola vez por clase de prueba;
        deben ser métodos \texttt{static} salvo que la clase esté anotada con
        \texttt{@TestInstance(Lifecycle.PER\_CLASS)}.
  \item \texttt{@DisplayName}: permite nombrar el caso de prueba con texto libre, incluidos
        caracteres especiales y espacios, lo que mejora la legibilidad de los reportes.
  \item \texttt{@Nested}: agrupa pruebas relacionadas en clases internas no estáticas; permite
        una jerarquía de contextos de prueba que refleja la lógica del dominio.
  \item \texttt{@ParameterizedTest} con \texttt{@ValueSource}, \texttt{@CsvSource} o
        \texttt{@MethodSource}: ejecuta el mismo caso de prueba con múltiples conjuntos de
        entrada, reduciendo la duplicación.
  \item \texttt{@Tag}: clasifica las pruebas para filtrado selectivo durante la ejecución.
  \item \texttt{@Disabled}: excluye temporalmente una prueba, con obligación de documentar
        la razón ---práctica que el autor recomienda imponer como parte de la rúbrica.
\end{itemize}

\subsection{Integración con Maven Surefire}
\label{subsec:maven-surefire}

La ejecución automática de la batería de pruebas se configura mediante el complemento
Maven Surefire. El fragmento del Listado~\ref{lst:surefire-config} muestra la configuración
mínima adoptada para el proyecto:

\begin{lstlisting}[style=javaStyle,
  caption={Configuración de Maven Surefire para JUnit~5 en \texttt{pom.xml}.},
  label={lst:surefire-config},
  language=XML]
<build>
  <plugins>
    <plugin>
      <groupId>org.apache.maven.plugins</groupId>
      <artifactId>maven-surefire-plugin</artifactId>
      <version>3.2.5</version>
      <configuration>
        <!-- Excluir pruebas de integracion marcadas con @Tag -->
        <excludedGroups>integracion</excludedGroups>
        <!-- Habilitar reporte en formato XML para SonarQube -->
        <reportFormat>xml</reportFormat>
        <!-- Mostrar nombres completos de pruebas fallidas -->
        <useFile>false</useFile>
      </configuration>
    </plugin>
  </plugins>
</build>
\end{lstlisting}

\noindent
Con esta configuración, el comando \texttt{mvn test} ejecuta todas las pruebas anotadas con
\texttt{@Test} y \texttt{@ParameterizedTest} que no pertenezcan al grupo \texttt{integracion}.
Los reportes XML generados en \texttt{target/surefire-reports/} son consumidos por SonarQube
para el análisis de cobertura y calidad (véase la sección~\ref{sec:cobertura}).

\subsection{Prueba unitaria de \texttt{Facultad}}
\label{subsec:prueba-facultad}

La clase \texttt{Facultad} es la raíz de la jerarquía organizacional del sistema.
Toda vez que el \acuerdo\ establece que toda unidad académica debe tener un código
institucional no nulo, que la Facultad contiene al menos una Escuela y que su director
debe ser designado formalmente, las pruebas han de verificar esas tres restricciones
de dominio con carácter prioritario. El Listado~\ref{lst:junit-facultad} muestra la
clase de prueba completa, incluyendo casos parametrizados y verificación de excepciones:

\begin{lstlisting}[style=javaStyle,
  caption={Clase de prueba \texttt{FacultadTest} con JUnit~5.},
  label={lst:junit-facultad}]
package co.edu.udistrital.estatuto.modelo;

import org.junit.jupiter.api.*;
import org.junit.jupiter.params.ParameterizedTest;
import org.junit.jupiter.params.provider.NullAndEmptySource;
import org.junit.jupiter.params.provider.ValueSource;
import static org.junit.jupiter.api.Assertions.*;
import java.time.LocalDate;

@DisplayName("Pruebas unitarias: Facultad")
class FacultadTest {

    private Facultad facultad;

    @BeforeEach
    void setUp() {
        facultad = new Facultad(
            "Facultad de Ingenieria",
            "FI-001",
            LocalDate.of(1948, 7, 20),
            "Decano de Prueba"
        );
    }

    @Test
    @DisplayName("getTipo() retorna 'Facultad'")
    void getTipoRetornaFacultad() {
        assertEquals("Facultad", facultad.getTipo(),
            "El tipo debe ser la cadena 'Facultad' segun el dominio");
    }

    @Test
    @DisplayName("Constructor: codigo nulo lanza IllegalArgumentException")
    void constructorCodigoNuloLanzaExcepcion() {
        assertThrows(IllegalArgumentException.class, () ->
            new Facultad(
                "Facultad de Tecnologia",
                null,                        // codigo nulo
                LocalDate.now(),
                "Decano"
            ),
            "Se esperaba IllegalArgumentException para codigo nulo"
        );
    }

    @ParameterizedTest
    @NullAndEmptySource
    @ValueSource(strings = {"  ", "\t", "\n"})
    @DisplayName("Constructor: codigo nulo o en blanco siempre lanza excepcion")
    void constructorCodigoInvalidoLanzaExcepcion(String codigoInvalido) {
        assertThrows(IllegalArgumentException.class, () ->
            new Facultad("Facultad X", codigoInvalido,
                         LocalDate.now(), "Director")
        );
    }

    @Test
    @DisplayName("addEscuela: agregar escuela incrementa la lista")
    void addEscuelaIncrementaLista() {
        Escuela esc = new Escuela(
            "Ingenieria de Sistemas",
            "IS-001",
            LocalDate.of(1994, 3, 1),
            "Director Sistemas",
            facultad
        );
        facultad.addEscuela(esc);
        assertEquals(1, facultad.getEscuelas().size());
        assertTrue(facultad.getEscuelas().contains(esc));
    }

    @Test
    @DisplayName("addEscuela: escuela nula lanza NullPointerException")
    void addEscuelaNulaLanzaExcepcion() {
        assertThrows(NullPointerException.class,
            () -> facultad.addEscuela(null));
    }

    @Test
    @DisplayName("getCodigo: el codigo se normaliza a mayusculas")
    void getCodigoNormalizadoMayusculas() {
        Facultad f = new Facultad("Tecnologia", "fi-002",
                                   LocalDate.now(), "Director");
        assertEquals("FI-002", f.getCodigo());
    }

    @Test
    @DisplayName("getFechaCreacion: devuelve la fecha pasada en el constructor")
    void getFechaCreacionCorrecta() {
        LocalDate esperada = LocalDate.of(1948, 7, 20);
        assertEquals(esperada, facultad.getFechaCreacion());
    }

    @Nested
    @DisplayName("Pruebas de igualdad y representacion")
    class IgualdadYRepresentacion {

        @Test
        @DisplayName("equals: dos facultades con mismo codigo son iguales")
        void equalsMismoCodigoSonIguales() {
            Facultad f2 = new Facultad("Nombre distinto", "FI-001",
                                        LocalDate.now(), "Otro director");
            assertEquals(facultad, f2,
                "La igualdad debe basarse en el codigo institucional");
        }

        @Test
        @DisplayName("toString: contiene nombre y codigo")
        void toStringContieneNombreYCodigo() {
            String repr = facultad.toString();
            assertTrue(repr.contains("FI-001") && repr.contains("Ingenieria"),
                "toString debe incluir nombre y codigo");
        }
    }
}
\end{lstlisting}

\subsection{Prueba unitaria de \texttt{Escuela}}
\label{subsec:prueba-escuela}

\begin{lstlisting}[style=javaStyle,
  caption={Clase de prueba \texttt{EscuelaTest} con JUnit~5.},
  label={lst:junit-escuela}]
package co.edu.udistrital.estatuto.modelo;

import org.junit.jupiter.api.*;
import static org.junit.jupiter.api.Assertions.*;
import java.time.LocalDate;

@DisplayName("Pruebas unitarias: Escuela")
class EscuelaTest {

    private Facultad facultadPadre;
    private Escuela escuela;

    @BeforeEach
    void setUp() {
        facultadPadre = new Facultad(
            "Facultad de Ingenieria", "FI-001",
            LocalDate.of(1948, 7, 20), "Decano"
        );
        escuela = new Escuela(
            "Ingenieria de Sistemas",
            "IS-001",
            LocalDate.of(1994, 3, 1),
            "Director Sistemas",
            facultadPadre
        );
    }

    @Test
    @DisplayName("getTipo() retorna 'Escuela'")
    void getTipoEsEscuela() {
        assertEquals("Escuela", escuela.getTipo());
    }

    @Test
    @DisplayName("getFacultad: la escuela conoce su facultad padre")
    void getfacultadCorrecta() {
        assertSame(facultadPadre, escuela.getFacultad(),
            "La escuela debe referenciar exactamente la facultad padre");
    }

    @Test
    @DisplayName("Constructor: facultad padre nula lanza NullPointerException")
    void constructorFacultadNulaLanzaExcepcion() {
        assertThrows(NullPointerException.class, () ->
            new Escuela("Telematica", "TL-001",
                        LocalDate.now(), "Director", null)
        );
    }

    @Test
    @DisplayName("addDocente: lista de docentes empieza vacia")
    void listaDocentesInicialmenteVacia() {
        assertTrue(escuela.getDocentes().isEmpty(),
            "Una escuela recien creada no tiene docentes asignados");
    }

    @Test
    @DisplayName("addDocente: agregar docente valido incrementa la lista")
    void addDocenteIncrementaLista() {
        Docente d = new Docente(
            "Ana", "Garcia",
            "ana.garcia@udistrital.edu.co",
            TipoVinculacion.PLANTA,
            escuela
        );
        escuela.addDocente(d);
        assertEquals(1, escuela.getDocentes().size());
    }

    @Test
    @DisplayName("No se puede agregar el mismo docente dos veces")
    void noSePuedeAgregarDocenteDuplicado() {
        Docente d = new Docente(
            "Luis", "Perez",
            "luis.perez@udistrital.edu.co",
            TipoVinculacion.OCASIONAL,
            escuela
        );
        escuela.addDocente(d);
        escuela.addDocente(d); // segundo intento
        assertEquals(1, escuela.getDocentes().size(),
            "La unicidad de adscripcion debe mantenerse");
    }
}
\end{lstlisting}

\subsection{Prueba unitaria de \texttt{Docente}}
\label{subsec:prueba-docente}

La verificación del enumerado \texttt{TipoVinculacion} es particularmente relevante desde
la perspectiva normativa: el \acuerdo\ reconoce exactamente cinco modalidades de vinculación
docente, y cualquier instancia con un valor fuera de ese conjunto constituiría una violación
del dominio. Dado que el enumerado en \java\ garantiza en tiempo de compilación que solo
se usen valores definidos, las pruebas se concentran en la lógica de negocio asociada
a cada tipo ---como el cálculo de la carga horaria máxima--- más que en la validez del valor
en sí.\footnote{Esta garantía en tiempo de compilación es una ventaja práctica de los
enumerados de \java\ frente a las representaciones numéricas o textuales que todavía
se encuentran en sistemas heredados. El autor ha detectado, en revisiones de código de
proyectos estudiantiles (2020--2025), que hasta el 40~\% de los estudiantes que no usan
enumerados representan el tipo de docente como un entero, lo que suscita errores de
dominio difíciles de rastrear.}

\begin{lstlisting}[style=javaStyle,
  caption={Clase de prueba \texttt{DocenteTest} con JUnit~5.},
  label={lst:junit-docente}]
package co.edu.udistrital.estatuto.modelo;

import org.junit.jupiter.api.*;
import org.junit.jupiter.params.ParameterizedTest;
import org.junit.jupiter.params.provider.EnumSource;
import static org.junit.jupiter.api.Assertions.*;
import java.time.LocalDate;

@DisplayName("Pruebas unitarias: Docente")
class DocenteTest {

    private Escuela escuelaPrueba;

    @BeforeEach
    void setUp() {
        Facultad f = new Facultad("Ingenieria", "FI-001",
                                   LocalDate.now(), "Decano");
        escuelaPrueba = new Escuela("Sistemas", "IS-001",
                                    LocalDate.now(), "Director", f);
    }

    @ParameterizedTest
    @EnumSource(TipoVinculacion.class)
    @DisplayName("Constructor: todo TipoVinculacion valido crea docente sin excepcion")
    void constructorAceptaTodosLosVinculos(TipoVinculacion tipo) {
        assertDoesNotThrow(() ->
            new Docente("Nombre", "Apellido",
                        "email@udistrital.edu.co", tipo, escuelaPrueba)
        );
    }

    @Test
    @DisplayName("getEmail: correo nulo lanza NullPointerException")
    void emailNuloLanzaExcepcion() {
        assertThrows(NullPointerException.class, () ->
            new Docente("Nombre", "Apellido", null,
                        TipoVinculacion.PLANTA, escuelaPrueba)
        );
    }

    @Test
    @DisplayName("getTipoVinculacion: retorna el tipo pasado en construccion")
    void getTipoVinculacionCorrecto() {
        Docente d = new Docente("Carlos", "Morales",
            "c.morales@udistrital.edu.co",
            TipoVinculacion.CATEDRA, escuelaPrueba);
        assertEquals(TipoVinculacion.CATEDRA, d.getTipoVinculacion());
    }

    @Test
    @DisplayName("esTiempoCompleto: solo docente de PLANTA es tiempo completo")
    void esTiempoCompletoSoloParaPlanta() {
        Docente planta = new Docente("A", "B", "a@ud.edu.co",
            TipoVinculacion.PLANTA, escuelaPrueba);
        Docente catedra = new Docente("C", "D", "c@ud.edu.co",
            TipoVinculacion.CATEDRA, escuelaPrueba);
        assertTrue(planta.esTiempoCompleto());
        assertFalse(catedra.esTiempoCompleto());
    }

    @Test
    @DisplayName("getNombreCompleto: concatena nombre y apellido correctamente")
    void getNombreCompletoFormato() {
        Docente d = new Docente("Maria", "Lopez",
            "m.lopez@ud.edu.co", TipoVinculacion.OCASIONAL, escuelaPrueba);
        assertEquals("Maria Lopez", d.getNombreCompleto());
    }
}
\end{lstlisting}

\subsection{Prueba unitaria de \texttt{Estudiante}}
\label{subsec:prueba-estudiante}

\begin{lstlisting}[style=javaStyle,
  caption={Clase de prueba \texttt{EstudianteTest} con JUnit~5.},
  label={lst:junit-estudiante}]
package co.edu.udistrital.estatuto.modelo;

import org.junit.jupiter.api.*;
import static org.junit.jupiter.api.Assertions.*;
import java.time.LocalDate;

@DisplayName("Pruebas unitarias: Estudiante")
class EstudianteTest {

    private Estudiante estudiante;

    @BeforeEach
    void setUp() {
        estudiante = new Estudiante(
            "Laura",
            "Ramirez",
            "20251001",          // codigo estudiantil
            1,                   // semestre actual
            EstadoEstudiante.ACTIVO
        );
    }

    @Test
    @DisplayName("getSemestre: semestre inicial es 1")
    void semestreInicialEsUno() {
        assertEquals(1, estudiante.getSemestre());
    }

    @Test
    @DisplayName("Constructor: semestre 0 lanza IllegalArgumentException")
    void semestreCeroLanzaExcepcion() {
        assertThrows(IllegalArgumentException.class, () ->
            new Estudiante("X", "Y", "20251002", 0, EstadoEstudiante.ACTIVO),
            "El semestre debe ser mayor que cero segun las reglas del dominio"
        );
    }

    @Test
    @DisplayName("Constructor: semestre negativo lanza IllegalArgumentException")
    void semestreNegativoLanzaExcepcion() {
        assertThrows(IllegalArgumentException.class, () ->
            new Estudiante("X", "Y", "20251003", -1, EstadoEstudiante.ACTIVO)
        );
    }

    @Test
    @DisplayName("getEstado: estado inicial es ACTIVO")
    void estadoInicialActivo() {
        assertEquals(EstadoEstudiante.ACTIVO, estudiante.getEstado());
    }

    @Test
    @DisplayName("avanzarSemestre: incrementa el semestre en 1")
    void avanzarSemestreIncrementa() {
        estudiante.avanzarSemestre();
        assertEquals(2, estudiante.getSemestre());
    }

    @Test
    @DisplayName("suspender: cambia estado a SUSPENDIDO")
    void suspenderCambiaEstado() {
        estudiante.suspender();
        assertEquals(EstadoEstudiante.SUSPENDIDO, estudiante.getEstado());
    }

    @Test
    @DisplayName("Constructor: codigo estudiantil nulo lanza NullPointerException")
    void codigoNuloLanzaExcepcion() {
        assertThrows(NullPointerException.class, () ->
            new Estudiante("A", "B", null, 1, EstadoEstudiante.ACTIVO)
        );
    }
}
\end{lstlisting}

Cuatro clases de prueba, diecinueve casos en total. La lectura de los listados anteriores
sugiere un patrón que conviene hacer explícito: cada método de prueba verifica exactamente
una condición, su nombre describe con precisión la situación y el resultado esperado, y
el código de preparación (\texttt{@BeforeEach}) se mantiene mínimo para evitar dependencias
entre pruebas. Esos tres hábitos ---granularidad, nomenclatura y aislamiento--- son,
a juicio del autor, más valiosos en la formación que cualquier herramienta concreta.


%% ─────────────────────────────────────────────
\section{Pruebas en \rust\ (\texttt{\#[test]}, \texttt{cargo test})}
\label{sec:pruebas-rust}
%% ─────────────────────────────────────────────

\subsection{Marco nativo de pruebas de \rust}
\label{subsec:rust-test-marco}

El mecanismo de pruebas de \rust\ se distingue del de \java\ en un aspecto estructural
determinante: está integrado en el compilador y en la cadena de herramientas estándar,
sin necesidad de dependencias externas. Tres atributos y dos macros conforman el vocabulario
básico:\footnote{Al momento de redactar este capítulo (marzo de 2026), \texttt{cargo test}
ejecuta los casos de prueba en hilos paralelos por defecto. Cuando las pruebas acceden a
recursos compartidos ---como archivos de configuración o bases de datos en memoria---,
es preciso añadir \texttt{-{}-test-threads=1} o sincronizar el acceso mediante
\texttt{std::sync::Mutex}. El autor ha verificado que este detalle es la causa más frecuente
de fallas intermitentes en proyectos de primer semestre que migran de \java\ a \rust.}

\begin{itemize}
  \item \texttt{\#[test]}: marca una función como caso de prueba. No requiere argumentos.
  \item \texttt{\#[cfg(test)]}: condiciona la compilación del bloque \texttt{mod tests}
        solo cuando se ejecuta \texttt{cargo test}; el código de prueba no se incluye
        en el binario de producción.
  \item \texttt{\#[should\_panic]}: indica que la prueba se considera exitosa únicamente
        si la función entra en pánico (\textit{panic}); equivale a \texttt{assertThrows}
        de JUnit~5.
  \item \texttt{assert!(\textit{cond})}: falla si la condición booleana es falsa.
  \item \texttt{assert\_eq!(\textit{a}, \textit{b})} / \texttt{assert\_ne!(\textit{a}, \textit{b})}:
        comparan por igualdad o desigualdad e imprimen ambos valores en caso de fallo,
        lo que facilita el diagnóstico.
\end{itemize}

\noindent
La sintaxis  \rust\ consiste en situar el módulo de pruebas al final
del mismo fichero que contiene el código de producción. Aunque es posible colocar las pruebas
en ficheros separados bajo el directorio \texttt{tests/}, la primera opción resulta más
habitual para pruebas unitarias, puesto que tiene acceso a los ítems privados del módulo
---algo que JUnit~5 logra solo mediante reflexión y modificación de visibilidad.

\subsection{Módulo de pruebas para \texttt{Facultad} y \texttt{Escuela}}
\label{subsec:rust-prueba-facultad-escuela}

\begin{lstlisting}[style=rustStyle,
  caption={Módulo de pruebas para los \textit{structs} \texttt{Facultad} y \texttt{Escuela} en \rust.},
  label={lst:rust-pruebas-facultad}]
#[cfg(test)]
mod tests {
    use super::*;
    use chrono::NaiveDate;

    // ── Funciones auxiliares (fixtures) ────────────────────────────────
    fn facultad_prueba() -> Facultad {
        Facultad::new(
            "Facultad de Ingenieria".to_string(),
            "FI-001".to_string(),
            NaiveDate::from_ymd_opt(1948, 7, 20).unwrap(),
            "Decano de Prueba".to_string(),
        ).expect("Facultad valida debe crearse sin error")
    }

    fn escuela_prueba(fac_id: u32) -> Escuela {
        Escuela::new(
            "Ingenieria de Sistemas".to_string(),
            "IS-001".to_string(),
            NaiveDate::from_ymd_opt(1994, 3, 1).unwrap(),
            "Director Sistemas".to_string(),
            fac_id,
        ).expect("Escuela valida debe crearse sin error")
    }

    // ── Pruebas de Facultad ─────────────────────────────────────────────

    #[test]
    fn facultad_tipo_correcto() {
        let f = facultad_prueba();
        assert_eq!(f.tipo(), "Facultad");
    }

    #[test]
    fn facultad_codigo_normalizado_mayusculas() {
        let f = Facultad::new(
            "Tecnologia".to_string(),
            "fi-002".to_string(),
            NaiveDate::from_ymd_opt(2000, 1, 1).unwrap(),
            "Director".to_string(),
        ).unwrap();
        assert_eq!(f.codigo(), "FI-002");
    }

    #[test]
    #[should_panic(expected = "codigo no puede estar vacio")]
    fn facultad_codigo_vacio_falla() {
        Facultad::new(
            "Nombre".to_string(),
            "".to_string(),
            NaiveDate::from_ymd_opt(2000, 1, 1).unwrap(),
            "Director".to_string(),
        ).unwrap();
    }

    #[test]
    fn facultad_sin_escuelas_al_crear() {
        let f = facultad_prueba();
        assert!(f.escuelas().is_empty(),
            "Una facultad nueva no debe tener escuelas");
    }

    #[test]
    fn facultad_agrega_escuela_correctamente() {
        let mut f = facultad_prueba();
        let esc = escuela_prueba(f.id());
        let n = f.agregar_escuela(esc);
        assert!(n.is_ok());
        assert_eq!(f.escuelas().len(), 1);
    }

    // ── Pruebas de Escuela ──────────────────────────────────────────────

    #[test]
    fn escuela_tipo_correcto() {
        let esc = escuela_prueba(1);
        assert_eq!(esc.tipo(), "Escuela");
    }

    #[test]
    fn escuela_conoce_su_facultad_padre() {
        let esc = escuela_prueba(42);
        assert_eq!(esc.facultad_id(), 42);
    }

    #[test]
    fn escuela_lista_docentes_inicialmente_vacia() {
        let esc = escuela_prueba(1);
        assert!(esc.docentes_ids().is_empty());
    }

    #[test]
    fn escuela_agrega_docente_sin_duplicar() {
        let mut esc = escuela_prueba(1);
        let _ = esc.agregar_docente_id(10);
        let _ = esc.agregar_docente_id(10); // duplicado
        assert_eq!(esc.docentes_ids().len(), 1,
            "La unicidad de adscripcion debe mantenerse en Rust tambien");
    }
}
\end{lstlisting}

\subsection{Módulo de pruebas para \texttt{Docente}}
\label{subsec:rust-prueba-docente}

\begin{lstlisting}[style=rustStyle,
  caption={Módulo de pruebas para el \textit{struct} \texttt{Docente} en \rust.},
  label={lst:rust-pruebas-docente}]
#[cfg(test)]
mod tests_docente {
    use super::*;

    fn docente_planta() -> Docente {
        Docente::new(
            "Carlos".to_string(),
            "Morales".to_string(),
            "c.morales@udistrital.edu.co".to_string(),
            TipoVinculacion::Planta,
            1, // escuela_id
        ).expect("Docente valido")
    }

    #[test]
    fn docente_planta_es_tiempo_completo() {
        let d = docente_planta();
        assert!(d.es_tiempo_completo());
    }

    #[test]
    fn docente_catedra_no_es_tiempo_completo() {
        let d = Docente::new(
            "Ana".to_string(), "Ruiz".to_string(),
            "a.ruiz@ud.edu.co".to_string(),
            TipoVinculacion::Catedra, 1,
        ).unwrap();
        assert!(!d.es_tiempo_completo());
    }

    #[test]
    fn docente_nombre_completo_formato_correcto() {
        let d = docente_planta();
        assert_eq!(d.nombre_completo(), "Carlos Morales");
    }

    #[test]
    #[should_panic(expected = "email no puede ser nulo")]
    fn docente_email_vacio_falla() {
        Docente::new(
            "X".to_string(), "Y".to_string(), "".to_string(),
            TipoVinculacion::Ocasional, 1,
        ).unwrap();
    }

    #[test]
    fn todos_los_tipos_de_vinculacion_crean_docente() {
        let tipos = [
            TipoVinculacion::Planta,
            TipoVinculacion::Ocasional,
            TipoVinculacion::Catedra,
            TipoVinculacion::Visitante,
            TipoVinculacion::Especial,
        ];
        for tipo in tipos {
            let resultado = Docente::new(
                "Nombre".to_string(), "Apellido".to_string(),
                "e@ud.edu.co".to_string(), tipo, 1,
            );
            assert!(resultado.is_ok(),
                "TipoVinculacion {:?} debe ser valido", tipo);
        }
    }
}
\end{lstlisting}

\subsection{Ejemplo de salida de \texttt{cargo test}}
\label{subsec:rust-cargo-test-output}

La ejecución de \texttt{cargo test -{}-{}-- -{}-nocapture} para el módulo anterior produce
una salida similar a la siguiente, donde cada prueba se identifica por su nombre completo:

\begin{lstlisting}[style=rustStyle,
  caption={Salida representativa de \texttt{cargo test} para el módulo de dominio.},
  label={lst:cargo-test-output}]
running 14 tests
test tests::facultad_tipo_correcto ... ok
test tests::facultad_codigo_normalizado_mayusculas ... ok
test tests::facultad_sin_escuelas_al_crear ... ok
test tests::facultad_agrega_escuela_correctamente ... ok
test tests::escuela_tipo_correcto ... ok
test tests::escuela_conoce_su_facultad_padre ... ok
test tests::escuela_lista_docentes_inicialmente_vacia ... ok
test tests::escuela_agrega_docente_sin_duplicar ... ok
test tests_docente::docente_planta_es_tiempo_completo ... ok
test tests_docente::docente_catedra_no_es_tiempo_completo ... ok
test tests_docente::docente_nombre_completo_formato_correcto ... ok
test tests_docente::todos_los_tipos_de_vinculacion_crean_docente ... ok
test tests::facultad_codigo_vacio_falla - should panic ... ok
test tests_docente::docente_email_vacio_falla - should panic ... ok

test result: ok. 14 passed; 0 failed; 0 ignored; 0 measured; 0 filtered out
\end{lstlisting}

\subsection{Contraste JUnit~5 / \rust\ nativo}
\label{subsec:contraste-junit-rust}

Ambos mecanismos verifican las mismas invariantes del dominio, pero difieren en la
disposición del código, en la integración con el compilador y en la granularidad del
diagnóstico de fallos. La tabla~\ref{tab:contraste-junit-rust} sintetiza las diferencias
más relevantes para un equipo de desarrollo universitario:

\begingroup
\small
\begin{longtable}{@{} L{3.5cm} L{5.5cm} L{5cm} @{} }
	% --- Título y encabezado de la primera página ---
	\caption{Contraste entre JUnit~5 y el mecanismo nativo de pruebas de \rust.} \label{tab:contraste-junit-rust} \\
	\toprule
	\textbf{Aspecto} & \textbf{JUnit~5 (\java)} & \textbf{\texttt{\#[test]} (\rust)} \\
	\midrule
	\endfirsthead
	
	% --- Encabezado para las páginas siguientes ---
	\multicolumn{3}{c}{{\bfseries \tablename\ \thetable{} -- Continuación}} \\
	\toprule
	\textbf{Aspecto} & \textbf{JUnit~5 (\java)} & \textbf{\texttt{\#[test]} (\rust)} \\
	\midrule
	\endhead
	
	% --- Pie de tabla para páginas intermedias ---
	\midrule
	\multicolumn{3}{r}{{Continúa en la siguiente página...}} \\
	\bottomrule
	\endfoot
	
	% --- Pie de tabla final ---
	\bottomrule
	\endlastfoot
	
	% --- Datos de la tabla ---
	Dependencia externa & 
	Sí (\texttt{junit-jupiter}). Requiere gestión vía Maven o Gradle & 
	No. El soporte de pruebas está integrado directamente en \texttt{rustc} y \texttt{cargo} \\
	\midrule
	
	Ubicación de pruebas & 
	Clases separadas en la jerarquía \texttt{src/test/java} & 
	En el mismo archivo fuente, típicamente en un módulo interno bajo \texttt{\#[cfg(test)]} \\
	\midrule
	
	Acceso a privados & 
	Limitado; requiere reflexión o cambiar visibilidad a \textit{package-private} & 
	Directo. El código de prueba en el mismo módulo tiene acceso total a miembros privados \\
	\midrule
	
	Pruebas parametrizadas & 
	Soporte nativo avanzado con \texttt{@ParameterizedTest} y \texttt{@ValueSource} & 
	Se logra mediante macros externas o bucles manuales sobre colecciones de datos \\
	\midrule
	
	Detección de errores & 
	Principalmente en tiempo de ejecución tras la compilación exitosa & 
	\texttt{cargo test} incluye la verificación de tipos y \textit{borrow checker} en el código de prueba \\
	\midrule
	
	Ejecución paralela & 
	Configurable mediante el motor Jupiter o plugins como Surefire & 
	Activada por defecto. Cada prueba corre en un hilo independiente gestionado por el \textit{test runner} \\
	\midrule
	
	Reportes e Integración & 
	Genera formatos estándar XML/HTML para herramientas de CI/CD (Jenkins, SonarQube) & 
	Salida en texto plano, TAP o JSON; integración nativa y transparente con cualquier CI \\
\end{longtable}
\endgroup


%% ─────────────────────────────────────────────
\section{Cobertura de código y métricas de calidad}
\label{sec:cobertura}
%% ─────────────────────────────────────────────

\subsection{JaCoCo para \java}
\label{subsec:jacoco}

JaCoCo (\textit{Java Code Coverage}) instrumenta el \textit{bytecode} de la JVM para
registrar qué instrucciones, ramas y líneas son ejecutadas durante la batería de pruebas.
La integración con Maven requiere añadir el complemento al \texttt{pom.xml} e indicar
los umbrales mínimos que deben cumplirse para que la construcción no falle:

\begin{lstlisting}[style=javaStyle,
  caption={Configuración de JaCoCo en \texttt{pom.xml} con umbrales de cobertura.},
  label={lst:jacoco-config},
  language=XML]
<plugin>
  <groupId>org.jacoco</groupId>
  <artifactId>jacoco-maven-plugin</artifactId>
  <version>0.8.11</version>
  <executions>
    <!-- Fase 1: preparar el agente de instrumentacion -->
    <execution>
      <id>prepare-agent</id>
      <goals><goal>prepare-agent</goal></goals>
    </execution>
    <!-- Fase 2: generar reporte HTML tras las pruebas -->
    <execution>
      <id>report</id>
      <phase>test</phase>
      <goals><goal>report</goal></goals>
    </execution>
    <!-- Fase 3: verificar umbrales minimos -->
    <execution>
      <id>check</id>
      <phase>verify</phase>
      <goals><goal>check</goal></goals>
      <configuration>
        <rules>
          <rule>
            <element>PACKAGE</element>
            <includes>
              <include>co/edu/udistrital/estatuto/modelo/**</include>
            </includes>
            <limits>
              <limit>
                <counter>LINE</counter>
                <value>COVEREDRATIO</value>
                <minimum>0.80</minimum>  <!-- 80% cobertura de linea -->
              </limit>
              <limit>
                <counter>BRANCH</counter>
                <value>COVEREDRATIO</value>
                <minimum>0.70</minimum>  <!-- 70% cobertura de rama -->
              </limit>
            </limits>
          </rule>
        </rules>
      </configuration>
    </execution>
  </executions>
</plugin>
\end{lstlisting}

\noindent
Los umbrales del 80~\% para cobertura de línea y el 70~\% para cobertura de rama se
establecen como valores mínimos para el paquete \texttt{modelo}. El reporte HTML generado
en \texttt{target/site/jacoco/} desglosa la cobertura por paquete, clase y método, e
identifica visualmente las ramas no ejecutadas mediante un código de color
(verde~= cubierta, amarillo~= parcialmente cubierta, rojo~= no cubierta).

Queda por determinar si esos umbrales son adecuados para un sistema con consecuencias
normativas como el que se desarrolla en este libro. La extrapolación exige cautela:
el 80~\% de cobertura de línea no implica que el 80~\% de las situaciones de uso real
estén probadas, toda vez que las entradas de usuario y las interacciones entre módulos
constituyen un espacio de comportamiento que las pruebas unitarias no agotan.

\subsection{\texttt{cargo-tarpaulin} para \rust}
\label{subsec:tarpaulin}

\texttt{cargo-tarpaulin} instrumenta el código \rust\ compilado para medir la cobertura
de línea. Su instalación y uso básico se muestran a continuación:

\begin{lstlisting}[style=rustStyle,
  caption={Instalación y ejecución de \texttt{cargo-tarpaulin}.},
  label={lst:tarpaulin-uso}]
# Instalacion (requiere Linux o macOS; en Windows usar WSL2)
cargo install cargo-tarpaulin

# Ejecucion basica: reporte en consola
cargo tarpaulin --verbose

# Reporte en formato XML (compatible con SonarQube y GitLab CI)
cargo tarpaulin --out Xml --output-dir target/coverage/

# Umbral: falla si la cobertura es inferior al 80%
cargo tarpaulin --fail-under 80

# Salida representativa
% cargo tarpaulin
Jun 2026 09:14:33 INFO  Running Tarpaulin
[...]
|| Tested/Total Lines:
|| src/modelo/facultad.rs: 42/51
|| src/modelo/escuela.rs: 38/44
|| src/modelo/docente.rs: 35/40
|| src/modelo/estudiante.rs: 29/34
||
82.14% coverage, 144/175 lines covered
\end{lstlisting}

\subsection{Métricas de calidad estructural}
\label{subsec:metricas-calidad}

La cobertura de pruebas es una condición necesaria pero insuficiente para garantizar la
calidad del código. Tres métricas estructurales complementan el análisis:\footnote{Las
métricas de calidad estructural ---complejidad ciclomática, acoplamiento entre objetos y
cohesión--- fueron formalizadas por McCabe (1976), Chidamber y Kemerer (1994) respectivamente.
El autor reconoce que no existe consenso sobre los umbrales óptimos para cada métrica en el
contexto de sistemas universitarios de mediana escala como el descrito en este libro.}

\begin{itemize}
  \item \textbf{Complejidad ciclomática (CC)}: cuenta el número de caminos independientes
        en un método. Una CC~$> 10$ en un método del dominio sugiere la necesidad de
        extracción de métodos auxiliares. Para el sistema del estatuto, se adopta CC~$\leq 8$
        como umbral en clases de dominio.
  \item \textbf{CBO (\textit{Coupling Between Objects})}: mide el número de clases con las
        que una clase dada está acoplada. Un CBO~$> 5$ en una clase del dominio es señal de
        responsabilidad excesiva; en controladores se tolera hasta CBO~$= 10$.
  \item \textbf{LCOM (\textit{Lack of Cohesion in Methods})}: valores cercanos a~1 indican
        baja cohesión (los métodos operan sobre conjuntos disjuntos de atributos). Para
        las clases del dominio se espera LCOM~$< 0{,}3$.
\end{itemize}

La tabla~\ref{tab:metricas-calidad} compara las herramientas disponibles para medir estas
métricas en ambos lenguajes:

\begingroup
\small
\begin{longtable}{@{} L{3.5cm} L{3.5cm} L{7cm} @{} }
	% --- Título y encabezado de la primera página ---
	\caption{Ecosistema de herramientas de calidad: \java\ vs.\ \rust.} \label{tab:metricas-calidad} \\
	\toprule
	\textbf{Herramienta} & \textbf{Lenguaje} & \textbf{Función principal} \\
	\midrule
	\endfirsthead
	
	% --- Encabezado para las páginas siguientes ---
	\multicolumn{3}{c}{{\bfseries \tablename\ \thetable{} -- Continuación}} \\
	\toprule
	\textbf{Herramienta} & \textbf{Lenguaje} & \textbf{Función principal} \\
	\midrule
	\endhead
	
	% --- Pie de tabla para páginas intermedias ---
	\midrule
	\multicolumn{3}{r}{{Continúa en la siguiente página...}} \\
	\bottomrule
	\endfoot
	
	% --- Pie de tabla final ---
	\bottomrule
	\endlastfoot
	
	% --- Herramientas para Java ---
	JaCoCo & \java & Medición de cobertura de código (línea, rama e instrucción) \\
	\midrule
	SonarQube & \java & Análisis de complejidad ciclomática (CC), acoplamiento (CBO), cohesión (LCOM), deuda técnica y duplicación \\
	\midrule
	SpotBugs & \java & Detección de defectos comunes de programación (\textit{NullPointerException}, problemas de concurrencia) \\
	\midrule
	Checkstyle & \java & Verificación automática de estilo y cumplimiento de convenciones de codificación \\
	\midrule \midrule
	
	% --- Herramientas para Rust ---
	\texttt{cargo-tarpaulin} & \rust & Análisis de cobertura de línea diseñado específicamente para el ecosistema de pruebas de Rust \\
	\midrule
	\texttt{clippy} & \rust & Análisis estático avanzado para detectar anti-patrones y ofrecer sugerencias de sintaxis \\
	\midrule
	\texttt{cargo-geiger} & \rust & Auditoría de seguridad para la detección y conteo de bloques de código \texttt{unsafe} en el grafo de dependencias \\
	\midrule
	\texttt{cargo-audit} & \rust & Identificación de vulnerabilidades conocidas en dependencias mediante la base de datos RUSTSEC \\
\end{longtable}
\endgroup

\noindent
La asimetría más notable entre ambos ecosistemas no reside en la cobertura ---ambos lenguajes
la miden con herramientas maduras--- sino en el análisis de calidad estructural: SonarQube
ofrece un modelo de calidad unificado con umbrales configurables y tableros de mando, mientras
que el ecosistema \rust\ distribuye esas responsabilidades entre múltiples herramientas sin un
punto de control centralizado. Para un proyecto universitario de alcance moderado, esa
diferencia podría ser irrelevante; en un contexto de integración continua con múltiples
contribuidores, resulta más determinante.


%% ─────────────────────────────────────────────
\section{Refactorización orientada a objetos}
\label{sec:refactorizacion}
%% ─────────────────────────────────────────────

La refactorización no altera el comportamiento externo del sistema; transforma la estructura
interna del código para hacerlo más comprensible, más mantenible y menos propenso a defectos.
El catálogo de Fowler \parencite{Fowler2018} proporciona un vocabulario compartido para
describir cada transformación, lo que facilita la comunicación entre miembros del equipo.
A continuación se presentan cinco tarjetas de refactorización aplicadas al código del
sistema del estatuto, cada una con el par antes/después en \java\ y en \rust.

\subsection{Tarjeta~1: \textit{Extract Method} (Extracción de método)}
\label{subsec:refact-extract-method}

\textbf{Motivación.} El constructor de \texttt{UnidadAcademica} concentra la validación de
sus argumentos en el cuerpo principal, lo que dificulta la reutilización de esa lógica y
aumenta la complejidad ciclomática del constructor. Se extrae la validación a un método
privado estático con nombre descriptivo.

\textbf{Mecánica.} Identificar el bloque de código que se desea extraer $\rightarrow$
crear un método privado con nombre que describe la intención $\rightarrow$ reemplazar
el bloque original por una llamada al nuevo método $\rightarrow$ ejecutar la batería de
pruebas para verificar que el comportamiento no cambia.

\begin{lstlisting}[style=javaStyle,
  caption={\textit{Extract Method} en \java: antes y después.},
  label={lst:refact-extract-java}]
// ── ANTES ─────────────────────────────────────────────────────────────────
protected UnidadAcademica(String nombre, String codigo,
                           LocalDate fechaCreacion, String director) {
    if (codigo == null || codigo.isBlank()) {
        throw new IllegalArgumentException(
            "El codigo no puede ser nulo ni vacio.");
    }
    if (nombre == null || nombre.isBlank()) {
        throw new IllegalArgumentException(
            "El nombre no puede ser nulo ni vacio.");
    }
    if (fechaCreacion == null) {
        throw new IllegalArgumentException(
            "La fecha de creacion no puede ser nula.");
    }
    this.codigo = codigo.trim().toUpperCase();
    this.nombre = nombre.trim();
    this.fechaCreacion = fechaCreacion;
    this.director = Objects.requireNonNull(director, "Director no puede ser nulo.");
}

// ── DESPUES ────────────────────────────────────────────────────────────────
protected UnidadAcademica(String nombre, String codigo,
                           LocalDate fechaCreacion, String director) {
    validarCamposObligatorios(nombre, codigo, fechaCreacion, director);
    this.codigo = codigo.trim().toUpperCase();
    this.nombre = nombre.trim();
    this.fechaCreacion = fechaCreacion;
    this.director = director;
}

private static void validarCamposObligatorios(String nombre, String codigo,
                                               LocalDate fecha, String director) {
    if (codigo == null || codigo.isBlank())
        throw new IllegalArgumentException("Codigo invalido.");
    if (nombre == null || nombre.isBlank())
        throw new IllegalArgumentException("Nombre invalido.");
    Objects.requireNonNull(fecha, "Fecha de creacion no puede ser nula.");
    Objects.requireNonNull(director, "Director no puede ser nulo.");
}
\end{lstlisting}

\begin{lstlisting}[style=rustStyle,
  caption={\textit{Extract Method} en \rust: extracción a función privada de validación.},
  label={lst:refact-extract-rust}]
// ── ANTES ─────────────────────────────────────────────────────────────────
impl UnidadBase {
    pub fn new(nombre: String, codigo: String,
               fecha_creacion: NaiveDate, director: String)
    -> Result<Self, String> {
        if codigo.is_empty() { return Err("codigo no puede estar vacio".to_string()); }
        if nombre.is_empty() { return Err("nombre no puede estar vacio".to_string()); }
        if director.is_empty() { return Err("director no puede estar vacio".to_string()); }
        Ok(UnidadBase { nombre, codigo: codigo.to_uppercase(), fecha_creacion, director })
    }
}

// ── DESPUES ────────────────────────────────────────────────────────────────
fn validar_campo_obligatorio(valor: &str, campo: &str) -> Result<(), String> {
    if valor.trim().is_empty() {
        Err(format!("{} no puede estar vacio", campo))
    } else {
        Ok(())
    }
}

impl UnidadBase {
    pub fn new(nombre: String, codigo: String,
               fecha_creacion: NaiveDate, director: String)
    -> Result<Self, String> {
        validar_campo_obligatorio(&codigo, "codigo")?;
        validar_campo_obligatorio(&nombre, "nombre")?;
        validar_campo_obligatorio(&director, "director")?;
        Ok(UnidadBase {
            nombre: nombre.trim().to_string(),
            codigo: codigo.trim().to_uppercase(),
            fecha_creacion,
            director,
        })
    }
}
\end{lstlisting}

\subsection{Tarjeta~2: \textit{Replace Conditional with Polymorphism} (Reemplazo de condicional por polimorfismo)}
\label{subsec:refact-polymorphism}

\textbf{Motivación.} La lógica de negocio relativa al tipo de docente aparece dispersa en
múltiples métodos como cadenas de \texttt{if (tipo == TipoVinculacion.PLANTA) \{...\}}.
Cada vez que se añade un nuevo tipo, todos esos métodos deben modificarse: violación
directa del principio abierto/cerrado (OCP). El reemplazo por polimorfismo centraliza
cada comportamiento en la subclase correspondiente.

\begin{lstlisting}[style=javaStyle,
  caption={\textit{Replace Conditional with Polymorphism} en \java: antes y después.},
  label={lst:refact-poly-java}]
// ── ANTES: logica distribuida con condicionales ─────────────────────────────
public int calcularHorasMaximas() {
    if (tipoVinculacion == TipoVinculacion.PLANTA) {
        return 40;
    } else if (tipoVinculacion == TipoVinculacion.OCASIONAL) {
        return 30;
    } else if (tipoVinculacion == TipoVinculacion.CATEDRA) {
        return 20;
    } else {
        return 16;
    }
}

// ── DESPUES: polimorfismo mediante clases especializadas ────────────────────
// Clase base abstracta
public abstract class Docente {
    public abstract int calcularHorasMaximas();
    // ... atributos y metodos comunes
}

public class DocentePlanta extends Docente {
    @Override public int calcularHorasMaximas() { return 40; }
}

public class DocenteOcasional extends Docente {
    @Override public int calcularHorasMaximas() { return 30; }
}

public class DocenteCatedra extends Docente {
    @Override public int calcularHorasMaximas() { return 20; }
}
\end{lstlisting}

\subsection{Tarjeta~3: \textit{Introduce Parameter Object} (Introducción de objeto parámetro)}
\label{subsec:refact-param-object}

\textbf{Motivación.} El constructor de \texttt{Docente} recibe múltiples parámetros
de tipo \texttt{String} que podrían intercambiarse por error. Agrupar los datos de
contacto en un objeto \texttt{DatosContacto} reduce esa fragilidad y mejora la
semántica del código.

\begin{lstlisting}[style=javaStyle,
  caption={\textit{Introduce Parameter Object} en \java: creación de \texttt{DatosContacto}.},
  label={lst:refact-param-obj-java}]
// ── ANTES ─────────────────────────────────────────────────────────────────
public Docente(String nombre, String apellido, String email,
               String telefono, String oficina,
               TipoVinculacion tipo, Escuela escuela) { ... }

// ── DESPUES ────────────────────────────────────────────────────────────────
public final class DatosContacto {
    private final String email;
    private final String telefono;
    private final String oficina;

    public DatosContacto(String email, String telefono, String oficina) {
        this.email = Objects.requireNonNull(email);
        this.telefono = telefono;  // puede ser nulo (docente sin telefono asignado)
        this.oficina = oficina;
    }
    // getters...
}

public Docente(String nombre, String apellido,
               DatosContacto contacto,
               TipoVinculacion tipo, Escuela escuela) { ... }
\end{lstlisting}

\subsection{Tarjeta~4: \textit{Extract Class} (Extracción de clase)}
\label{subsec:refact-extract-class}

\textbf{Motivación.} La clase \texttt{Persona} acumula responsabilidades relacionadas con
la identidad personal (nombre, documento) y con el contacto institucional (email, teléfono,
oficina). Esas dos responsabilidades tienen razones distintas para cambiar ---violación del
principio de responsabilidad única (SRP)---. Se extrae \texttt{DatosContacto} como clase
independiente.

\begin{lstlisting}[style=rustStyle,
  caption={\textit{Extract Class} en \rust: separación de \texttt{DatosContacto}.},
  label={lst:refact-extract-class-rust}]
// ── ANTES: struct Persona con demasiadas responsabilidades ─────────────────
#[derive(Debug, Clone)]
pub struct Persona {
    pub nombre: String,
    pub apellido: String,
    pub documento: String,
    pub email: String,       // responsabilidad de contacto
    pub telefono: String,    // responsabilidad de contacto
    pub oficina: String,     // responsabilidad de contacto
}

// ── DESPUES: separacion de responsabilidades ───────────────────────────────
#[derive(Debug, Clone)]
pub struct IdentidadPersonal {
    pub nombre: String,
    pub apellido: String,
    pub documento: String,
}

#[derive(Debug, Clone)]
pub struct DatosContacto {
    pub email: String,
    pub telefono: Option<String>,
    pub oficina: Option<String>,
}

#[derive(Debug, Clone)]
pub struct Persona {
    pub identidad: IdentidadPersonal,
    pub contacto: DatosContacto,
}
\end{lstlisting}

\subsection{Tarjeta~5: \textit{Move Method} (Traslado de método)}
\label{subsec:refact-move-method}

\textbf{Motivación.} El método \texttt{calcularPromedioDocentes()} se encuentra en la clase
\texttt{FacultadControlador}, pero opera exclusivamente sobre datos de \texttt{Facultad}.
Al estar en el controlador, la lógica de negocio queda imbricada con la infraestructura,
dificultando las pruebas unitarias y violando el DIP. Se traslada al objeto de dominio.

\begin{lstlisting}[style=javaStyle,
  caption={\textit{Move Method} en \java: traslado de lógica al dominio.},
  label={lst:refact-move-method-java}]
// ── ANTES: metodo de negocio en el controlador ─────────────────────────────
public class FacultadControlador {
    public double calcularPromedioDocentes(Facultad facultad) {
        if (facultad.getEscuelas().isEmpty()) return 0.0;
        return facultad.getEscuelas().stream()
            .mapToInt(esc -> esc.getDocentes().size())
            .average()
            .orElse(0.0);
    }
}

// ── DESPUES: logica trasladada al objeto de dominio ────────────────────────
public class Facultad extends UnidadAcademica {
    // ...
    public double calcularPromedioDocentes() {
        if (escuelas.isEmpty()) return 0.0;
        return escuelas.stream()
            .mapToInt(esc -> esc.getDocentes().size())
            .average()
            .orElse(0.0);
    }
}

public class FacultadControlador {
    // El controlador delega; ya no contiene logica de negocio
    public double obtenerPromedioDocentes(String codigoFacultad) {
        Facultad f = dao.buscarPorCodigo(codigoFacultad)
                        .orElseThrow(EntidadNoEncontradaException::new);
        return f.calcularPromedioDocentes(); // delega al dominio
    }
}
\end{lstlisting}


%% ─────────────────────────────────────────────
\section{Principios SOLID aplicados al sistema del Estatuto}
\label{sec:solid}
%% ─────────────────────────────────────────────

Los principios SOLID \parencite{Martin2017} no son reglas absolutas, sino heurísticas que
orientan el diseño hacia sistemas más mantenibles, extensibles y comprensibles. Su aplicación
al dominio institucional del estatuto presenta matices que vale la pena examinar con cierto
detenimiento, toda vez que el sistema tiene un alcance bien acotado ---el entramado operativo
de una universidad pública colombiana--- y restricciones normativas explícitas.

\begingroup
\small
\begin{longtable}{@{} L{1.2cm} L{3cm} L{5.5cm} L{4.3cm} @{} }
	% --- Título y encabezado de la primera página ---
	\caption{Principios SOLID: evaluación en el sistema del Estatuto UD.} \label{tab:solid-evaluacion} \\
	\toprule
	\textbf{Sigla} & \textbf{Principio} & \textbf{Ejemplo en el sistema} & \textbf{Clase / interfaz} \\
	\midrule
	\endfirsthead
	
	% --- Encabezado para las páginas siguientes ---
	\multicolumn{4}{c}{{\bfseries \tablename\ \thetable{} -- Continuación}} \\
	\toprule
	\textbf{Sigla} & \textbf{Principio} & \textbf{Ejemplo en el sistema} & \textbf{Clase / interfaz} \\
	\midrule
	\endhead
	
	% --- Pie de tabla para páginas intermedias ---
	\midrule
	\multicolumn{4}{r}{{Continúa en la siguiente página...}} \\
	\bottomrule
	\endfoot
	
	% --- Pie de tabla final ---
	\bottomrule
	\endlastfoot
	
	% --- Datos de la tabla ---
	SRP & Responsabilidad única &
	La entidad \texttt{UnidadAcademica} cambia solo si la estructura organizacional del dominio varía &
	\texttt{UnidadAcademica}, \texttt{DatosContacto} \\
	\midrule
	
	OCP & Abierto/cerrado &
	Posibilidad de añadir un \texttt{Laboratorio} o \texttt{CABA} sin modificar la lógica de creación existente &
	\texttt{UnidadAcademicaFactory}, \texttt{Laboratorio} \\
	\midrule
	
	LSP & Sustitución de Liskov &
	\texttt{Facultad} y \texttt{Escuela} son intercambiables en cualquier módulo que procese una \texttt{UnidadAcademica} &
	Jerarquía de \texttt{UnidadAcademica} \\
	\midrule
	
	ISP & Segregación de interfaces &
	Dividir \texttt{GenericDAO<T>} para evitar que entidades de solo lectura implementen métodos de escritura &
	\texttt{GenericDAO}, \texttt{ReadOnlyDAO} \\
	\midrule
	
	DIP & Inversión de dependencias &
	\texttt{FacultadControlador} depende de la abstracción \texttt{GenericDAO}, no de una implementación JSON específica &
	\texttt{GenericDAO<Facultad>}, \texttt{FacultadJsonDAO} \\
\end{longtable}
\endgroup

\subsection{SRP: Responsabilidad única}
\label{subsec:solid-srp}

La clase \texttt{UnidadAcademica} encapsula datos del dominio organizacional. Su única razón
para cambiar es una modificación en la estructura de las unidades académicas definida por el
\acuerdo. Si se le añadiera lógica de presentación ---como la generación de un PDF con el
organigrama--- o de persistencia ---como la serialización a JSON---, tendría múltiples razones
para cambiar y el SRP quedaría comprometido. La separación entre \texttt{modelo},
\texttt{controlador} y \texttt{persistencia} en la arquitectura del Capítulo~\ref{cap:java}
responde precisamente a ese imperativo.

Pero la aplicación del SRP no resulta siempre inequívoca. Queda abierta la pregunta de si
\texttt{UnidadAcademica} debería incluir el método \texttt{calcularPromedioDocentes()}
---trasladado allí en la Tarjeta~5 de refactorización--- o si ese cálculo pertenece a una
capa de servicio separada. El autor no tiene una respuesta definitiva; la ubicación más
adecuada depende del contexto de evolución del sistema.

\subsection{OCP: Abierto/cerrado}
\label{subsec:solid-ocp}

El \acuerdo\ define cuatro tipos de unidades académicas: Facultad, Escuela, Centro e Instituto.
Si en el futuro se añade un quinto tipo ---pongamos, \texttt{Laboratorio}--- el sistema debe
poder acomodarlo sin modificar las clases existentes. El patrón Factory Method, empleado en
\texttt{UnidadAcademicaFactory}, lo posibilita:

\begin{lstlisting}[style=javaStyle,
  caption={Factory Method extensible para unidades académicas.},
  label={lst:ocp-factory}]
public class UnidadAcademicaFactory {
    public static UnidadAcademica crear(String tipo, String nombre,
                                        String codigo, LocalDate fecha,
                                        String director, Object... extras) {
        return switch (tipo.toUpperCase()) {
            case "FACULTAD"  -> new Facultad(nombre, codigo, fecha, director);
            case "ESCUELA"   -> new Escuela(nombre, codigo, fecha, director,
                                            (Facultad) extras[0]);
            case "CENTRO"    -> new Centro(nombre, codigo, fecha, director);
            case "INSTITUTO" -> new Instituto(nombre, codigo, fecha, director);
            // Para agregar Laboratorio, solo se aniade un nuevo case:
            // case "LABORATORIO" -> new Laboratorio(nombre, codigo, fecha, director);
            default -> throw new IllegalArgumentException(
                "Tipo de unidad desconocido: " + tipo);
        };
    }
}
\end{lstlisting}

\subsection{LSP: Sustitución de Liskov}
\label{subsec:solid-lsp}

El principio de Liskov establece que toda subclase debe poder emplearse donde se espera
la superclase, sin que el cliente note la diferencia. Para la jerarquía del estatuto esto
implica que \texttt{Facultad}, \texttt{Escuela}, \texttt{Centro} e \texttt{Instituto}
deben responder correctamente al mensaje \texttt{getTipo()} y respetar las postcondiciones
de \texttt{UnidadAcademica} ---en particular, que el código no es nulo y está normalizado
a mayúsculas. Las pruebas parametrizadas con \texttt{@ParameterizedTest} son el mecanismo
más adecuado para verificar el LSP de forma sistemática:

\begin{lstlisting}[style=javaStyle,
  caption={Verificación del LSP mediante prueba parametrizada.},
  label={lst:lsp-prueba}]
@ParameterizedTest
@MethodSource("proveedorDeUnidades")
@DisplayName("LSP: toda UnidadAcademica devuelve codigo no nulo en mayusculas")
void todaUnidadRetornaCodigoNormatizado(UnidadAcademica unidad) {
    assertNotNull(unidad.getCodigo());
    assertEquals(unidad.getCodigo(), unidad.getCodigo().toUpperCase());
    assertFalse(unidad.getCodigo().isBlank());
}

static Stream<UnidadAcademica> proveedorDeUnidades() {
    Facultad f = new Facultad("FI", "FI-001", LocalDate.now(), "Decano");
    return Stream.of(
        f,
        new Escuela("IS", "IS-001", LocalDate.now(), "Dir", f),
        new Centro("CI", "CI-001", LocalDate.now(), "Dir"),
        new Instituto("II", "II-001", LocalDate.now(), "Dir")
    );
}
\end{lstlisting}

\subsection{ISP: Segregación de interfaces}
\label{subsec:solid-isp}

La interfaz \texttt{GenericDAO<T>} define operaciones de lectura y escritura:
\texttt{guardar}, \texttt{buscarPorId}, \texttt{buscarTodos}, \texttt{actualizar}
y \texttt{eliminar}. Para entidades que solo se consultan ---como los artículos del
estatuto, que son datos de referencia--- implementar \texttt{guardar} y \texttt{eliminar}
constituiría una violación del ISP: la implementación forzaría a lanzar
\texttt{UnsupportedOperationException}, señal inequívoca de interfaz demasiado amplia.
La solución consiste en segregar \texttt{GenericDAO<T>} en dos interfaces:

\begin{lstlisting}[style=javaStyle,
  caption={Segregación de \texttt{GenericDAO} conforme al ISP.},
  label={lst:isp-segregacion}]
// Interfaz de solo lectura
public interface ReadOnlyDAO<T, ID> {
    Optional<T> buscarPorId(ID id);
    List<T> buscarTodos();
}

// Interfaz de escritura (extiende la de lectura)
public interface GenericDAO<T, ID> extends ReadOnlyDAO<T, ID> {
    void guardar(T entidad);
    void actualizar(T entidad);
    void eliminar(ID id);
}

// Las entidades de referencia usan solo ReadOnlyDAO
public class ArticuloEstatutoDAO implements ReadOnlyDAO<ArticuloEstatuto, Integer> { ... }

// Las entidades del dominio usan GenericDAO completo
public class FacultadJsonDAO implements GenericDAO<Facultad, String> { ... }
\end{lstlisting}

\subsection{DIP: Inversión de dependencias}
\label{subsec:solid-dip}

El principio de inversión de dependencias requiere que las clases de alto nivel ---los
controladores--- dependan de abstracciones (interfaces), no de implementaciones concretas.
En el sistema del estatuto, \texttt{FacultadControlador} recibe su dependencia de acceso
a datos mediante inyección en el constructor, lo que permite sustituir \texttt{FacultadJsonDAO}
por \texttt{FacultadMemoriaDAO} en las pruebas sin modificar el controlador:

\begin{lstlisting}[style=javaStyle,
  caption={DIP: controlador con dependencia inyectada sobre interfaz.},
  label={lst:dip-controlador}]
public class FacultadControlador {
    // Depende de la abstraccion, no de la implementacion concreta
    private final GenericDAO<Facultad, String> dao;

    // Inyeccion de dependencias por constructor
    public FacultadControlador(GenericDAO<Facultad, String> dao) {
        this.dao = Objects.requireNonNull(dao);
    }

    public Facultad buscarFacultad(String codigo) {
        return dao.buscarPorId(codigo)
                  .orElseThrow(() -> new EntidadNoEncontradaException(
                      "Facultad con codigo " + codigo + " no encontrada."));
    }
}

// En pruebas: inyectar un DAO en memoria, no el DAO JSON real
class FacultadControladorTest {
    @Test
    void buscarFacultadExistenteRetornaEntidad() {
        GenericDAO<Facultad, String> daoMock = new FacultadMemoriaDAO();
        daoMock.guardar(new Facultad("FI", "FI-001", LocalDate.now(), "D"));
        FacultadControlador ctrl = new FacultadControlador(daoMock);
        assertNotNull(ctrl.buscarFacultad("FI-001"));
    }
}
\end{lstlisting}


%% ─────────────────────────────────────────────
\section{Documentación del código (Javadoc, rustdoc)}
\label{sec:documentacion}
%% ─────────────────────────────────────────────

La documentación del código es, en el contexto del sistema del estatuto, un requisito de
dominio además de una buena práctica de ingeniería: toda entidad del sistema tiene un
correlato normativo en el \acuerdo, y esa correspondencia debe quedar explícita en los
comentarios de la clase o el \textit{struct}.\footnote{El autor considera que la documentación
más valiosa en código orientado a dominio no es la que describe \emph{cómo} hace algo un
método, sino \emph{por qué} el dominio exige ese comportamiento. La referencia explícita
al artículo del estatuto aplicable, tal como se muestra en los listados siguientes, es un
ejemplo de esa práctica.}

\subsection{Javadoc: documentación de \texttt{Facultad}}
\label{subsec:javadoc-facultad}

\begin{lstlisting}[style=javaStyle,
  caption={Documentación Javadoc completa para la clase \texttt{Facultad}.},
  label={lst:javadoc-facultad}]
package co.edu.udistrital.estatuto.modelo;

import java.time.LocalDate;
import java.util.*;

/**
 * Representa una Facultad de la Universidad Distrital Francisco José de Caldas.
 *
 * <p>Conforme al artículo~5 del Acuerdo 004 de 2025 (Estatuto General), una Facultad
 * es la unidad académico-administrativa de mayor jerarquía en la estructura de la
 * universidad. Agrupa Escuelas, Centros e Institutos bajo la dirección de un Decano
 * designado por el Consejo Superior Universitario.</p>
 *
 * <p><strong>Invariantes del dominio:</strong></p>
 * <ul>
 *   <li>El código institucional no puede ser nulo ni vacío; se normaliza a mayúsculas.</li>
 *   <li>Una Facultad puede tener cero o más Escuelas asociadas.</li>
 *   <li>El Decano es un dato obligatorio en el momento de creación.</li>
 * </ul>
 *
 * @author Roberto Albeiro Pava-Díaz
 * @version 2.0
 * @since Estatuto General 2025 (Acuerdo 004 de 2025)
 * @see UnidadAcademica
 * @see Escuela
 */
public class Facultad extends UnidadAcademica {

    /** Lista de Escuelas adscritas a esta Facultad. Puede estar vacía. */
    private final List<Escuela> escuelas;

    /**
     * Construye una Facultad con sus datos básicos.
     *
     * @param nombre        nombre oficial de la Facultad (p.~ej., "Facultad de Ingeniería")
     * @param codigo        código institucional único (se normaliza a mayúsculas; no nulo ni vacío)
     * @param fechaCreacion fecha de creación o reconocimiento oficial
     * @param director      nombre del Decano vigente (no nulo)
     * @throws IllegalArgumentException si el código es nulo o vacío
     * @throws NullPointerException     si nombre, fechaCreacion o director son nulos
     */
    public Facultad(String nombre, String codigo,
                    LocalDate fechaCreacion, String director) {
        super(nombre, codigo, fechaCreacion, director);
        this.escuelas = new ArrayList<>();
    }

    /**
     * {@inheritDoc}
     *
     * @return la cadena {@code "Facultad"}, identificador del tipo de unidad
     */
    @Override
    public String getTipo() { return "Facultad"; }

    /**
     * Asocia una Escuela a esta Facultad.
     *
     * <p>La unicidad de la asociación se garantiza por comparación de código institucional.
     * Si la Escuela ya está asociada, el método retorna sin modificar la lista.</p>
     *
     * @param escuela la Escuela que se desea asociar (no nula)
     * @throws NullPointerException si {@code escuela} es {@code null}
     */
    public void addEscuela(Escuela escuela) {
        Objects.requireNonNull(escuela, "La escuela no puede ser nula.");
        if (!escuelas.contains(escuela)) escuelas.add(escuela);
    }

    /**
     * Retorna una vista no modificable de las Escuelas asociadas.
     *
     * @return lista inmutable de Escuelas; nunca {@code null}
     */
    public List<Escuela> getEscuelas() {
        return Collections.unmodifiableList(escuelas);
    }

    /**
     * Calcula el promedio de docentes por Escuela.
     *
     * @return promedio aritmético de docentes por Escuela, o {@code 0.0} si
     *         no hay Escuelas asociadas
     */
    public double calcularPromedioDocentes() {
        if (escuelas.isEmpty()) return 0.0;
        return escuelas.stream()
                       .mapToInt(e -> e.getDocentes().size())
                       .average()
                       .orElse(0.0);
    }
}
\end{lstlisting}

\subsection{rustdoc: documentación de \texttt{Facultad}}
\label{subsec:rustdoc-facultad}

\begin{lstlisting}[style=rustStyle,
  caption={Documentación \texttt{rustdoc} completa para el \textit{struct} \texttt{Facultad}.},
  label={lst:rustdoc-facultad}]
/// Representa una Facultad de la Universidad Distrital Francisco José de Caldas.
///
/// Conforme al artículo 5 del Acuerdo 004 de 2025 (Estatuto General), una Facultad
/// es la unidad académico-administrativa de mayor jerarquía en la estructura de la
/// universidad. Agrupa Escuelas bajo la dirección de un Decano.
///
/// # Invariantes del dominio
///
/// - El `codigo` no puede estar vacío; se normaliza a mayúsculas en la construcción.
/// - `escuelas_ids` puede ser un vector vacío, pero nunca `None`.
/// - El `decano` es un campo obligatorio.
///
/// # Errores
///
/// [`Facultad::new`] retorna `Err(String)` si `codigo` o `nombre` están vacíos.
///
/// # Ejemplo
///
/// ```rust
/// use chrono::NaiveDate;
/// use modelo::Facultad;
///
/// let f = Facultad::new(
///     "Facultad de Ingeniería".to_string(),
///     "FI-001".to_string(),
///     NaiveDate::from_ymd_opt(1948, 7, 20).unwrap(),
///     "Decano Vigente".to_string(),
/// ).expect("Facultad válida");
///
/// assert_eq!(f.tipo(), "Facultad");
/// assert_eq!(f.codigo(), "FI-001");
/// ```
#[derive(Debug, Clone, serde::Serialize, serde::Deserialize)]
pub struct Facultad {
    id: u32,
    nombre: String,
    codigo: String,
    fecha_creacion: chrono::NaiveDate,
    director: String,
    escuelas_ids: Vec<u32>,
}

impl Facultad {
    /// Construye una nueva Facultad validando los campos obligatorios.
    ///
    /// # Parámetros
    ///
    /// - `nombre`: nombre oficial de la Facultad.
    /// - `codigo`: código institucional único (se normaliza a mayúsculas).
    /// - `fecha_creacion`: fecha de creación o reconocimiento oficial.
    /// - `director`: nombre del Decano vigente.
    ///
    /// # Errores
    ///
    /// Retorna `Err` si `codigo` o `nombre` son cadenas vacías o solo contienen espacios.
    pub fn new(nombre: String, codigo: String,
               fecha_creacion: chrono::NaiveDate,
               director: String) -> Result<Self, String> {
        validar_campo_obligatorio(&codigo, "codigo")?;
        validar_campo_obligatorio(&nombre, "nombre")?;
        Ok(Facultad {
            id: generar_id(),
            nombre: nombre.trim().to_string(),
            codigo: codigo.trim().to_uppercase(),
            fecha_creacion,
            director,
            escuelas_ids: Vec::new(),
        })
    }

    /// Retorna el tipo de unidad académica.
    ///
    /// Siempre retorna `"Facultad"`.
    pub fn tipo(&self) -> &str { "Facultad" }

    /// Retorna el código institucional normalizado a mayúsculas.
    pub fn codigo(&self) -> &str { &self.codigo }

    /// Asocia una Escuela a esta Facultad por su identificador.
    ///
    /// Si el `id` ya estaba registrado, no se duplica.
    pub fn agregar_escuela(&mut self, escuela: Escuela) -> Result<(), String> {
        if !self.escuelas_ids.contains(&escuela.id()) {
            self.escuelas_ids.push(escuela.id());
        }
        Ok(())
    }
}
\end{lstlisting}

\noindent
Para generar la documentación HTML de \rust\ se emplea el comando \texttt{cargo doc -{}-open},
que compila la documentación de todos los módulos del proyecto y la abre en el navegador.
Los ejemplos incrustados en los bloques de código de la documentación (\textit{doc-tests})
se ejecutan como pruebas ordinarias durante \texttt{cargo test}, lo que garantiza que los
ejemplos de la documentación estén siempre sincronizados con el código real ---una
ventaja que Javadoc no ofrece de forma nativa.


%% ─────────────────────────────────────────────
\section{Comparativa \java\ vs.\ \rust: ecosistema de pruebas y calidad}
\label{sec:java-rust-pruebas}
%% ─────────────────────────────────────────────

Pese a que los principios de prueba unitaria son independientes del lenguaje ---un caso de
prueba verifica una condición, de forma aislada, con un resultado determinista---, el
ecosistema que rodea esas pruebas difiere de manera notable entre \java\ y \rust. Esas
diferencias inciden sobre la curva de aprendizaje, la integración en flujos de trabajo
de desarrollo continuo y la accesibilidad para estudiantes de pregrado con experiencia limitada.

\begingroup
\small
\begin{longtable}{@{} L{3cm} L{4.5cm} L{4.5cm} L{1.5cm} @{}}
	% --- Título y encabezado de la primera página ---
	\caption{Ecosistema de pruebas y calidad: \java\ vs.\ \rust\ (comparativa).} \label{tab:ecosistema-pruebas} \\
	\toprule
	\textbf{Dimensión} & \textbf{\java\ + JUnit~5} & \textbf{\rust\ + \texttt{\#[test]}} & \textbf{Ventaja} \\
	\midrule
	\endfirsthead
	
	% --- Encabezado para las páginas siguientes ---
	\multicolumn{4}{c}{{\bfseries \tablename\ \thetable{} -- Continuación}} \\
	\toprule
	\textbf{Dimensión} & \textbf{\java\ + JUnit~5} & \textbf{\rust\ + \texttt{\#[test]}} & \textbf{Ventaja} \\
	\midrule
	\endhead
	
	% --- Pie de tabla para páginas intermedias ---
	\midrule
	\multicolumn{4}{r}{{Continúa en la siguiente página...}} \\
	\bottomrule
	\endfoot
	
	% --- Pie de tabla final ---
	\bottomrule
	\endlastfoot
	
	% --- Datos de la tabla ---
	\textit{Framework} & 
	JUnit~5 (dependencia externa vía Maven/Gradle) & 
	Nativo en el compilador \texttt{rustc} y \texttt{cargo} & 
	\rust \\
	\midrule
	
	Pruebas de dobles & 
	Mockito (estándar de facto, muy maduro en el ecosistema) & 
	\texttt{mockall} (dependencia externa en \texttt{Cargo.toml}) & 
	\java \\
	\midrule
	
	Cobertura & 
	JaCoCo + complemento de construcción Maven & 
	\texttt{cargo-tarpaulin} (análisis de líneas ejecutadas) & 
	Paridad \\
	\midrule
	
	Análisis estático & 
	SpotBugs, Checkstyle, SonarQube (muy robustos) & 
	\texttt{clippy}, \texttt{cargo-audit} (integrados y rápidos) & 
	\java \\
	\midrule
	
	Configuración & 
	Moderada (gestión de dependencias en \texttt{pom.xml}) & 
	Mínima (soporte de primer nivel en \texttt{Cargo.toml}) & 
	\rust \\
	\midrule
	
	Integración CI/CD & 
	Alta soporte en GitHub Actions, Jenkins y GitLab & 
	Alta mediante \texttt{cargo test} en cualquier entorno CI & 
	Paridad \\
	\midrule
	
	Diagnóstico & 
	Moderada (basada en \textit{stack traces} de la JVM) & 
	Alta (mensajes detallados del compilador y \textit{panic}) & 
	\rust \\
	\midrule
	
	Curva de aprendizaje & 
	Moderada (conceptos conocidos de la JVM y Maven) & 
	Pronunciada (\textit{ownership} y \textit{borrow checker}) & 
	\java \\
	\midrule
	
	Madurez & 
	Muy alta (más de 25 años de evolución estable) & 
	Creciente y acelerada (desde la versión 1.0 en 2015) & 
	\java \\
\end{longtable}
\endgroup

\noindent
La conclusión más llamativa de la comparativa no es cuál ecosistema es «mejor» en términos
absolutos, sino cuál resulta más adecuado en función del perfil del equipo y del momento del
aprendizaje. En las primeras semanas del proyecto, cuando los estudiantes están familiarizando-
se con los conceptos de orientación a objetos, el ecosistema de \java\ ---más verboso pero
más transparente en su configuración--- reduce la fricción cognitiva. En etapas avanzadas,
cuando la corrección del código es prioritaria, el \textit{borrow checker} de \rust\ ofrece
garantías en tiempo de compilación que ningún \textit{framework} de pruebas puede reemplazar.

\subsection{Integración en flujos de CI/CD}
\label{subsec:cicd}

Ambos lenguajes integran con naturalidad en flujos de integración continua basados en
GitHub Actions. El Listado~\ref{lst:github-actions} muestra un flujo mínimo que ejecuta
pruebas en los dos lenguajes:

\begin{lstlisting}[style=javaStyle,
  caption={Flujo de GitHub Actions para pruebas en \java\ y \rust\ simultáneamente.},
  label={lst:github-actions},
  language=XML]
name: CI — Pruebas Java y Rust

on: [push, pull_request]

jobs:
  java:
    runs-on: ubuntu-latest
    steps:
      - uses: actions/checkout@v4
      - uses: actions/setup-java@v4
        with: { java-version: '21', distribution: 'temurin' }
      - run: cd proyecto_java && mvn verify

  rust:
    runs-on: ubuntu-latest
    steps:
      - uses: actions/checkout@v4
      - uses: dtolnay/rust-toolchain@stable
      - run: cd proyecto_rust && cargo test --all
      - run: cargo clippy -- -D warnings
\end{lstlisting}


%% ─────────────────────────────────────────────────────
%% SPRINT 10
%% ─────────────────────────────────────────────────────
\section*{Sprint 10: Pruebas unitarias y cobertura}
\label{sprint:10}
\addcontentsline{toc}{section}{Sprint 10: Pruebas unitarias y cobertura}

\begin{tcolorbox}[sprintBox, title={Sprint 10 --- Pruebas unitarias y cobertura}]
  \textbf{Duración estimada:} 1 semana (semanas~15--16 del semestre)\\
  \textbf{Objetivo:} Implementar la batería de pruebas unitarias para los módulos
  organizacional y académico del sistema en \java\ y en \rust, alcanzar al menos
  80~\% de cobertura de línea en las clases del dominio, configurar JaCoCo y
  \texttt{cargo-tarpaulin}, y documentar los casos de prueba en la matriz de trazabilidad
  del proyecto.\\
  \textbf{Prerequisito:} Código funcional de los Sprints~7 y~8 (Capítulos~\ref{cap:java}
  y~\ref{cap:rust}), con al menos los módulos \texttt{modelo} y \texttt{controlador}
  implementados y compilando sin errores.
\end{tcolorbox}

\bigskip

\subsection*{Tabla del Sprint 10}

\begin{longtable}{L{3.5cm} L{4.5cm} L{4.5cm}}
  \sprintcabecera
  \endhead
  %
  Clase \texttt{FacultadTest} en JUnit~5: mínimo 6 casos de prueba que cubran
  constructor válido, código nulo, código vacío, addEscuela, igualdad y toString
  &
  Verificar las invariantes del dominio de la entidad organizacional de mayor
  jerarquía, conforme al artículo~5 del \acuerdo
  &
  Todas las pruebas pasan con \texttt{mvn test}; la clase \texttt{FacultadTest}
  cubre al menos 6 casos con \texttt{@DisplayName} descriptivos; reporte JaCoCo
  indica $\geq 80$~\% de cobertura de línea en \texttt{Facultad} \\
  \midrule
  %
  Clases \texttt{EscuelaTest}, \texttt{DocenteTest} y \texttt{EstudianteTest}
  en JUnit~5: mínimo 5 casos de prueba cada una, incluyendo al menos 2 casos
  con \texttt{assertThrows}
  &
  Verificar las invariantes de dominio de las entidades académicas y la
  validez del enumerado \texttt{TipoVinculacion}
  &
  Todas las pruebas de las tres clases pasan; al menos un caso
  \texttt{@ParameterizedTest} con \texttt{@EnumSource} verifica todos los
  valores de \texttt{TipoVinculacion}; la cobertura de rama de \texttt{Docente}
  supera el 70~\% \\
  \midrule
  %
  Módulo de pruebas \rust\ (\texttt{\#[cfg(test)]}) para \texttt{Facultad},
  \texttt{Escuela} y \texttt{Docente}: mínimo 4 casos por \textit{struct},
  incluyendo al menos 1 con \texttt{\#[should\_panic]}
  &
  Aplicar el mecanismo nativo de pruebas de \rust\ al dominio del estatuto,
  verificando las mismas invariantes que en \java\ con las herramientas de sintaxis del
  lenguaje
  &
  \texttt{cargo test} pasa sin fallos; al menos 1 \textit{doc-test} en la
  documentación de \texttt{Facultad} se ejecuta y pasa; salida de
  \texttt{cargo tarpaulin} indica $\geq 75$~\% de cobertura \\
  \midrule
  %
  Reporte de cobertura JaCoCo integrado en Maven (\texttt{mvn verify}) con
  umbrales configurados: 80~\% línea y 70~\% rama para el paquete
  \texttt{modelo}; construcción falla si no se alcanzan los umbrales
  &
  Institucionalizar la cobertura como umbral de calidad verificable
  automáticamente en cada ciclo de integración
  &
  El archivo \texttt{pom.xml} incluye la configuración de JaCoCo con los
  umbrales documentados; \texttt{mvn verify} falla con mensaje claro si la
  cobertura es insuficiente; el reporte HTML se genera en
  \texttt{target/site/jacoco/} \\
  \midrule
  %
  Matriz de trazabilidad prueba--invariante--artículo: tabla que relaciona
  cada caso de prueba con la invariante que verifica y el artículo del
  \acuerdo\ del que se deriva
  &
  Documentar la correspondencia entre los casos de prueba automatizados y las
  restricciones normativas del estatuto universitario
  &
  La tabla cubre al menos 15 casos de prueba; cada fila identifica la prueba,
  la invariante verificada y el artículo del \acuerdo; el documento se
  adjunta al repositorio en \texttt{docs/trazabilidad\_pruebas.md} \\
  \bottomrule
  \caption{Entregables, características y rúbrica del Sprint~10.}
  \label{tab:sprint10}
\end{longtable}

\subsection*{Actividad de reflexión: software generativo y pruebas unitarias}
\label{act:ia-sprint10}

\begin{tcolorbox}[actividadIA,
  title={Actividad de reflexión (Sprint~10) ---
         Software generativo y pruebas unitarias}]

\textbf{Instrucción general.}
El estudiante diseña y escribe de manera individual las pruebas unitarias de la clase
\texttt{Facultad}, identificando previamente las invariantes del dominio que deben
verificarse y consultando el \acuerdo. A continuación, solicita a un asistente generativo
(por ejemplo, GitHub Copilot, ChatGPT o Claude) que genere las mismas pruebas para la misma
clase, proporcionándole únicamente el código fuente de \texttt{Facultad} como contexto.
El estudiante compara ambos conjuntos de pruebas y elabora un informe de análisis crítico.

\bigskip
\textbf{Preguntas orientadoras:}
\begin{enumerate}[label=\arabic*.]
  \item ¿El asistente generativo produce pruebas que verifican invariantes del dominio
        institucional (código no nulo, normalización a mayúsculas, unicidad de adscripción)
        o se limita a pruebas triviales de \textit{getters} y \textit{setters}?
  \item ¿Las pruebas generadas son independientes entre sí, o presentan dependencias de
        orden de ejecución que podrían generar fallos intermitentes?
  \item ¿El asistente asigna nombres descriptivos que comunican la intención del caso de
        prueba (\texttt{codigoNuloLanzaIllegalArgumentException}) o usa nombres genéricos
        (\texttt{test1}, \texttt{testConstructor})?
  \item ¿Hay pruebas que el asistente omitió y que el estudiante identificó a partir de
        la lectura del \acuerdo? ¿A qué atribuye esa omisión?
  \item ¿Las pruebas generadas por el asistente compilarían y pasarían con el código de
        producción actual, o requieren ajustes antes de poder ejecutarse?
\end{enumerate}

\bigskip
\textbf{Producto esperado:}
\begin{itemize}
  \item Clase \texttt{FacultadTest} escrita por el estudiante (en el repositorio).
  \item Transcripción o captura de pantalla de la sesión con el asistente generativo.
  \item Informe comparativo de 400--600 palabras que responda las preguntas orientadoras
        con evidencia concreta de los dos conjuntos de pruebas.
\end{itemize}

\bigskip
\textbf{Rúbrica:}
\begin{itemize}
  \item \textbf{Producto} (60~\%): cobertura de invariantes del dominio, independencia
        de los casos de prueba, nomenclatura descriptiva, ausencia de código duplicado.
  \item \textbf{Proceso} (40~\%): identificación de pruebas faltantes en la propuesta
        del asistente, análisis crítico fundamentado, conclusión sobre la utilidad
        y las limitaciones del asistente generativo para esta tarea.
\end{itemize}

\end{tcolorbox}


%% ─────────────────────────────────────────────────────
%% SPRINT 11
%% ─────────────────────────────────────────────────────
\section*{Sprint 11: Refactorización y principios SOLID}
\label{sprint:11}
\addcontentsline{toc}{section}{Sprint 11: Refactorización y principios SOLID}

\begin{tcolorbox}[sprintBox, title={Sprint 11 --- Refactorización y principios SOLID}]
  \textbf{Duración estimada:} 1 semana (semanas~17--18 del semestre)\\
  \textbf{Objetivo:} Refactorizar el código de los módulos organizacional y académico
  aplicando al menos tres principios SOLID, documentar las refactorizaciones con tarjetas
  antes/después siguiendo el formato del catálogo de Fowler \parencite{Fowler2018}, y
  verificar que la batería de pruebas del Sprint~10 sigue pasando íntegramente tras
  cada refactorización.\\
  \textbf{Prerequisito:} Batería de pruebas del Sprint~10 completa y en verde
  (\texttt{mvn test} y \texttt{cargo test} pasan sin fallos).
\end{tcolorbox}

\bigskip

\subsection*{Tabla del Sprint 11}

\begin{longtable}{L{3.5cm} L{4.5cm} L{4.5cm}}
  \sprintcabecera
  \endhead
  %
  Informe de refactorización con mínimo 5 tarjetas (nombre / motivación / mecánica
  / antes / después) sustentadas en principios SOLID; cada tarjeta incluye fragmentos
  de código en \java\ \emph{y} en \rust
  &
  Aplicar técnicas de refactorización al código del dominio institucional con
  fundamentación en el catálogo de Fowler y los principios SOLID; demostrar
  que la técnica es independiente del lenguaje
  &
  Las tarjetas identifican correctamente el problema de diseño, justifican la
  refactorización con un principio SOLID nombrado explícitamente, y documentan el
  cambio con fragmentos de código antes/después; el código resultante compila
  sin errores en ambos lenguajes \\
  \midrule
  %
  Batería de pruebas del Sprint~10 ejecutada sin cambios y pasando tras cada
  refactorización; evidencia: captura de reporte de ejecución antes y después
  &
  Garantizar que la refactorización no introduce regresiones; demostrar que
  una batería de pruebas sólida es el andamiaje que hace posible la refactorización
  segura
  &
  Todas las pruebas del Sprint~10 pasan tras la refactorización; el reporte de
  cobertura JaCoCo no decrece con respecto al obtenido antes de la refactorización;
  ambas capturas (antes y después) se adjuntan al informe \\
  \midrule
  %
  Evaluación de los cinco principios SOLID frente al código del sistema: tabla
  que identifica, para cada principio, si se cumple, se cumple parcialmente o
  se viola, con justificación y referencia al código
  &
  Desarrollar la capacidad de evaluar críticamente el diseño de un sistema real
  frente a criterios de calidad establecidos
  &
  La tabla es completa (cinco principios), la justificación es específica
  (referencia a clases e interfaces concretas del sistema), y las valoraciones
  son coherentes con el código entregado en el repositorio \\
  \midrule
  %
  Documentación Javadoc y rustdoc actualizada tras la refactorización: todas
  las clases y \textit{structs} del módulo \texttt{modelo} tienen documentación
  en el formato descrito en la sección~\ref{sec:documentacion}
  &
  Mantener la documentación sincronizada con el código, incluyendo la referencia
  al artículo del estatuto aplicable en cada entidad del dominio
  &
  \texttt{mvn javadoc:javadoc} genera documentación sin advertencias;
  \texttt{cargo doc} genera documentación sin errores; al menos el 80~\%
  de los métodos públicos tienen documentación con \texttt{@param}, \texttt{@return}
  y \texttt{@throws} (Javadoc) o sección \textbf{Errores} (rustdoc) \\
  \midrule
  %
  Análisis de complejidad ciclomática: al menos un método del módulo \texttt{modelo}
  con CC~$>$~5 antes de la refactorización; evidencia de que la refactorización
  redujo la CC al umbral adoptado ($\leq 8$ para clases de dominio)
  &
  Aplicar métricas cuantitativas de calidad para fundamentar decisiones de
  refactorización; no refactorizar por intuición, sino por evidencia medible
  &
  El informe incluye la CC del método antes y después de la refactorización;
  la medición se obtiene de SonarQube (para \java) o de \texttt{cargo clippy}
  con la opción \texttt{-{}-warn complexity} (para \rust) \\
  \bottomrule
  \caption{Entregables, características y rúbrica del Sprint~11.}
  \label{tab:sprint11}
\end{longtable}

\subsection*{Actividad de reflexión: software generativo y refactorización}
\label{act:ia-sprint11}

\begin{tcolorbox}[actividadIA,
  title={Actividad de reflexión (Sprint~11) ---
         Software generativo y refactorización}]

\textbf{Instrucción general.}
El estudiante identifica de manera individual una violación del principio abierto/cerrado
(OCP) o del principio de responsabilidad única (SRP) en el código desarrollado en los
sprints anteriores, y la refactoriza aplicando la técnica del catálogo de Fowler que
considere más adecuada. A continuación, describe el problema al asistente generativo
---incluyendo el fragmento de código con el problema--- y solicita una propuesta de
refactorización alternativa. El estudiante compara su solución con la del asistente.

\bigskip
\textbf{Preguntas orientadoras:}
\begin{enumerate}[label=\arabic*.]
  \item ¿La refactorización propuesta por el asistente aplica la técnica correcta del
        catálogo de Fowler para el tipo de problema identificado, o propone una técnica
        diferente? ¿Cuál es más adecuada y por qué?
  \item ¿La propuesta del asistente preserva el comportamiento externo del código,
        verificable mediante las pruebas del Sprint~10, o introduce cambios de semántica
        que haría fallar alguna prueba?
  \item ¿El asistente identifica en el código problemas de diseño adicionales que el
        estudiante no había observado? ¿Son esos problemas legítimos o son ruido?
  \item ¿La solución del asistente introduce dependencias nuevas no justificadas,
        incrementa la complejidad ciclomática o viola algún otro principio SOLID?
  \item En retrospectiva: ¿hubiese sido posible evitar el problema de diseño identificado
        si se hubieran aplicado los principios SOLID desde el Sprint~7?
\end{enumerate}

\bigskip
\textbf{Producto esperado:}
\begin{itemize}
  \item Tarjeta de refactorización propia (nombre / motivación / mecánica / antes / después).
  \item Propuesta del asistente generativo (transcripción o captura de pantalla).
  \item Informe comparativo de 400--600 palabras que responda las cinco preguntas
        orientadoras con evidencia concreta del código.
  \item Evidencia de que las pruebas del Sprint~10 siguen pasando tras la refactorización.
\end{itemize}

\bigskip
\textbf{Rúbrica:}
\begin{itemize}
  \item \textbf{Producto} (60~\%): corrección de la refactorización propia, preservación
        del comportamiento, aplicación correcta del principio SOLID, calidad del código
        resultante.
  \item \textbf{Proceso} (40~\%): evaluación crítica de la propuesta del asistente,
        identificación de fortalezas y debilidades, capacidad de argumentar la decisión
        de diseño adoptada.
\end{itemize}

\end{tcolorbox}

\bigskip

%% ─────────────────────────────────────────────
%% CIERRE DEL CAPÍTULO
%% ─────────────────────────────────────────────

\noindent
Los dos sprints de este capítulo cierran el ciclo de desarrollo principal del sistema
del estatuto.\footnote{La integración entre las herramientas de análisis estático
---SonarQube para \java, \texttt{clippy} para \rust--- y los flujos de CI/CD descritos
en la sección~\ref{subsec:cicd} puede resultar compleja para equipos sin experiencia
previa. El autor recomienda introducir SonarQube en una instalación local (Docker) antes
de intentar la integración con el servidor de CI, de manera que los estudiantes comprendan
la interfaz de usuario antes de interpretar los reportes en el \textit{pipeline}.} La batería
de pruebas construida en el Sprint~10 no es un producto estático: debe crecer junto con el
código, de modo que cada nueva funcionalidad incorpore su correspondiente caso de prueba.
La refactorización del Sprint~11, a su vez, no es un ejercicio de cosmética del código,
sino la demostración práctica de que el diseño orientado a objetos es un proceso iterativo
que se afina con la experiencia y con la retroalimentación que solo las pruebas automatizadas
pueden proporcionar de forma fiable.

Persiste, sin embargo, una duda que el autor no desea soslayar: la refactorización guiada
por principios SOLID puede conducir a una proliferación de clases e interfaces que incrementa
la complejidad estructural sin beneficio proporcional en mantenibilidad. La regla práctica
adoptada en la cátedra ---no crear una abstracción hasta que existan al menos dos implementaciones
concretas--- no es universal, pero sí ha probado ser útil en proyectos de alcance universitario
donde el tiempo de desarrollo es el recurso más escaso. El Capítulo~\ref{cap:rust} complementa
esta perspectiva desde el punto de vista del sistema de tipos de \rust, que impone disciplina
de diseño no mediante principios abstractos, sino mediante restricciones del compilador.
